\section{Sparse Inference and Pruning}
\label{sec:efficiency}

The implicit factor graph is sparse: each node has only a few neighbors,
and most conditional relationships are zero or near-zero. If the sparse
factor graph structure can be identified --- which SAE features matter for
a given task, which AND connections are active, which OR gates are relevant
--- the rest can be skipped. This is structured pruning with a theoretical
justification.

\subsection*{Task-Specific Factor Graphs}

Different tasks activate different subgraphs of the implicit factor graph.
This is the theoretical foundation for mixture-of-experts and sparse
attention: they are approximations to task-specific factor graph routing,
and the BP framework gives a principled account of what they are
approximating.

\subsection*{Distillation as Factor Graph Compression}

Knowledge distillation can be understood as factor graph compression:
finding a smaller factor graph that approximates the posteriors of the
larger one. The BP framework gives a principled objective for distillation:
minimize KL divergence between the posteriors of the teacher and student
factor graphs, not just between their output distributions.

\subsection*{Epistemic Status}

\begin{center}
\begin{tabular}{ll}
\toprule
Claim & Status \\
\midrule
Implicit factor graph is sparse & Mechanistic interpretation \\
Task-specific subgraph activation motivates sparse computation & Mechanistic interpretation \\
MoE approximates task-specific factor graph routing & Mechanistic interpretation \\
Distillation as factor graph compression & Architectural proposal \\
\bottomrule
\end{tabular}
\end{center}